% Options for packages loaded elsewhere
\PassOptionsToPackage{unicode}{hyperref}
\PassOptionsToPackage{hyphens}{url}
\documentclass[
]{article}
\usepackage{xcolor}
\usepackage{amsmath,amssymb}
\setcounter{secnumdepth}{5}
\usepackage{iftex}
\ifPDFTeX
  \usepackage[T1]{fontenc}
  \usepackage[utf8]{inputenc}
  \usepackage{textcomp} % provide euro and other symbols
\else % if luatex or xetex
  \usepackage{unicode-math} % this also loads fontspec
  \defaultfontfeatures{Scale=MatchLowercase}
  \defaultfontfeatures[\rmfamily]{Ligatures=TeX,Scale=1}
\fi
% \usepackage{lmodern}
\ifPDFTeX\else
  % xetex/luatex font selection
\fi
% Use upquote if available, for straight quotes in verbatim environments
\IfFileExists{upquote.sty}{\usepackage{upquote}}{}
\IfFileExists{microtype.sty}{% use microtype if available
  \usepackage[]{microtype}
  \UseMicrotypeSet[protrusion]{basicmath} % disable protrusion for tt fonts
}{}
\makeatletter
\@ifundefined{KOMAClassName}{% if non-KOMA class
  \IfFileExists{parskip.sty}{%
    \usepackage{parskip}
  }{% else
    \setlength{\parindent}{0pt}
    \setlength{\parskip}{6pt plus 2pt minus 1pt}}
}{% if KOMA class
  \KOMAoptions{parskip=half}}
\makeatother
\usepackage{longtable,booktabs,array}
\newcounter{none} % for unnumbered tables
\usepackage{calc} % for calculating minipage widths
% Correct order of tables after \paragraph or \subparagraph
\usepackage{etoolbox}
\makeatletter
\patchcmd\longtable{\par}{\if@noskipsec\mbox{}\fi\par}{}{}
\makeatother
% Allow footnotes in longtable head/foot
\IfFileExists{footnotehyper.sty}{\usepackage{footnotehyper}}{\usepackage{footnote}}
\makesavenoteenv{longtable}
\setlength{\emergencystretch}{3em} % prevent overfull lines
\providecommand{\tightlist}{%
  \setlength{\itemsep}{0pt}\setlength{\parskip}{0pt}}
\usepackage[]{natbib}
\bibliographystyle{plainnat}
\makeatletter
\@ifpackageloaded{subcaption}{}{\usepackage{subcaption}}
\@ifpackageloaded{caption}{}{\usepackage{caption}}
\captionsetup[subfigure]{margin=0.5em}
\AtBeginDocument{%
\renewcommand*\figurename{Figure}
\renewcommand*\tablename{Table}
}
\AtBeginDocument{%
\renewcommand*\listfigurename{List of Figures}
\renewcommand*\listtablename{List of Tables}
}
\newcounter{pandoccrossref@subfigures@footnote@counter}
\newenvironment{pandoccrossrefsubfigures}{%
\setcounter{pandoccrossref@subfigures@footnote@counter}{0}
\begin{figure}\centering%
\gdef\global@pandoccrossref@subfigures@footnotes{}%
\DeclareRobustCommand{\footnote}[1]{\footnotemark%
\stepcounter{pandoccrossref@subfigures@footnote@counter}%
\ifx\global@pandoccrossref@subfigures@footnotes\empty%
\gdef\global@pandoccrossref@subfigures@footnotes{{##1}}%
\else%
\g@addto@macro\global@pandoccrossref@subfigures@footnotes{, {##1}}%
\fi}}%
{\end{figure}%
\addtocounter{footnote}{-\value{pandoccrossref@subfigures@footnote@counter}}
\@for\f:=\global@pandoccrossref@subfigures@footnotes\do{\stepcounter{footnote}\footnotetext{\f}}%
\gdef\global@pandoccrossref@subfigures@footnotes{}}
\@ifpackageloaded{float}{}{\usepackage{float}}
\floatstyle{ruled}
\@ifundefined{c@chapter}{\newfloat{codelisting}{h}{lop}}{\newfloat{codelisting}{h}{lop}[chapter]}
\floatname{codelisting}{Listing}
\newcommand*\listoflistings{\listof{codelisting}{List of Listings}}
\makeatother
\usepackage{bookmark}
\IfFileExists{xurl.sty}{\usepackage{xurl}}{} % add URL line breaks if available
\urlstyle{same}
% fallback for those not using the hyperref driver hyperxmp:
\makeatletter
\@ifundefined{xmpquote}{\newcommand{\xmpquote}[1]{#1}}{}
\makeatother
\hypersetup{
  colorlinks=true,linkcolor=red,urlcolor=red,citecolor=red,anchorcolor=red,filecolor=red,
  pdfcreator={LaTeX via pandoc}}

\author{}
\date{}

\IfFileExists{seqsplit.sty}{\usepackage{seqsplit}}{\newcommand{\seqsplit}[1]{#1}}
\protected\def\breakseq#1{\seqsplit{#1}}
\protected\def\breaktt#1{\begingroup\ttfamily\seqsplit{#1}\endgroup}
\title{The Shape Between: A Full-Page Illustrated Storybook Template}
\author{Daniel Ari Friedman\\\footnotesize{Active Inference Institute}\\\footnotesize{\texttt{daniel@activeinference.institute}}\\\footnotesize{\href{https://orcid.org/0000-0001-6232-9096}{ORCID: 0000-0001-6232-9096}} \\ \footnotesize{DOI: \href{https://doi.org/10.5281/zenodo.21176000}{10.5281/zenodo.21176000}}}
\date{\today}

\usepackage[margin=0.75in]{geometry}
\usepackage{graphicx}

\begin{document}

\begin{titlepage}
\centering
\vspace*{2cm}
{\LARGE\bfseries The Shape Between: A Full-Page Illustrated Storybook Template\par}
\vspace{0.75em}
{\large A geometric fable of belonging, bracing, and reciprocal form\par}
\vfill
\makeatletter
{\@author\par}
\vspace{1em}
{\@date\par}
\makeatother
\vfill
\end{titlepage}
\thispagestyle{empty}
\tableofcontents
\newpage


\section{Abstract}\label{abstract}

\breaktt{template\_storybook} demonstrates a public, standalone
picture-book workflow in the research template repository. The bundled
project renders a deterministic fourteen-page storybook subtitled
\emph{A geometric fable of belonging, bracing, and reciprocal form}: a
clear cover, a publication-and-acknowledgements page, and twelve story
pages in which a child tetrahedron raised by cubes meets a child cube
raised by tetrahedra. The story uses a large reciprocal symbol, a
tetrahedron-inside-cube stability spread, a shadow-projection lesson, a
tensegrity lantern, and a vector garden to frame belonging without
sameness: square and triangle families remain distinct while the space
between them becomes navigable.

The project separates story data, rendering logic, and orchestration.
Story text, characters, page order, palettes, and overlay choices live
in \breaktt{content/story.yaml}; character generation and illustration
live in \texttt{src/storybook/}; one script renders the cover, each page
script renders a numbered page, and a final script assembles the PDF.
This makes the exemplar a forkable pattern for full-page creative
artifacts rather than standard manuscript-centered figure generation.

\newpage

\section{Introduction}\label{introduction}

Most template exemplars in this repository produce analytical reports,
manuscripts, newspapers, textbooks, or structured research artifacts. A
storybook stresses a different surface: the designed page is the
artifact. The page must be a complete visual composition, text must
remain readable, and the pipeline must still expose deterministic source
data, tests, and generated outputs.

The bundled story is intentionally symbolic. Tessa is a tetrahedron in a
family of cubes. Ciro is a cube in a tetrahedral family. Their meeting
is arranged around a yin-yang-like valley so that each family contains a
trace of the other without erasing its own geometry. The point is not
mathematical instruction; it is a compact visual grammar for reciprocal
belonging.

\newpage

\section{Architecture}\label{architecture}

The storybook uses the same Layer-2 project contract as the other public
templates:

\begin{itemize}
\tightlist
\item
  \breaktt{content/story.yaml} owns cast records and page records.
\item
  \breaktt{src/storybook/characters.py} validates shape characters.
\item
  \breaktt{src/storybook/story.py} loads the content file into typed
  records.
\item
  \breaktt{src/storybook/illustration.py} renders the full-page scenes.
\item
  \breaktt{src/storybook/rendering.py} assembles PNG pages into the PDF
  and writes the manifest.
\item
  \texttt{scripts/} contains the thin Stage-02 entry points.
\end{itemize}

The page scripts are deliberately repetitive. Their job is visible
orchestration: cover, page 1, page 2, and so on. The repetition keeps
the creative page surface easy to inspect while preventing scripts from
becoming hidden rendering engines.

\newpage

\section{Story and Graphics}\label{story-and-graphics}

Every page is rendered as a full-page PNG. Some pages place text
directly over the illustration with shadowed lettering; others use a
translucent box when the scene is visually dense. The YAML
\texttt{overlay\_box} value controls that choice per page.

\subsection{Methods}\label{methods}

The renderer uses simple geometric primitives: isometric cubes, shaded
tetrahedra, family clusters, a curved reciprocal symbol, deterministic
star fields, projected shadows, suspended tensegrity struts, and radial
vector gardens. The stability page borrows from the Synergetics
intuition that triangulated structure braces square space: a tetrahedron
is drawn through four alternating cube corners so the cube becomes
steadier without ceasing to be a cube \citep{fuller1975synergetics}.
Later pages extend that grammar through projection, push-pull balance,
and directional growth. This gives forks a readable baseline. Replacing
the story does not require editing Python; replacing the visual grammar
does.

\newpage

\section{Reproducibility}\label{reproducibility}

The project avoids external image services and live APIs. Page images
are procedurally generated from checked-in content, and the final PDF is
assembled from those generated images. Tests check that the story loads,
the two central characters have opposite family shapes, a page renders
at the configured dimensions, the PDF page count matches the content
file, and at least one page script can run against a temporary project
root.

The generated manifest records the cast, page list, and render paths.
That manifest is the compact evidence surface for the primary artifact.

\newpage

\section{References}\label{references}

The bundled story is illustrative fiction. The Synergetics reference
grounds the tetrahedron-inside-cube design motif used in the stability
page and the later tensegrity/vector motifs.



\clearpage

\bibliography{references}
\end{document}
